% A stub of Chapter 2, so that notation-check has a second chapter to compare.
% Invented sample: every name, number and reference is made up.
\chapter{Background: one queue, several desks}
\label{ch:background}

This chapter sets out the notation that the rest of the thesis uses.
Borrowers arrive at a return desk at random, at an average rate of $\lambda$ returns per hour.
Each desk serves returns at an average rate of $\mu$ returns per hour.
The library opens $N$ desks, each staffed by one volunteer.

The utilisation of the desks is the share of time that they are busy:
\begin{equation}
  \rho = \frac{\lambda}{N\mu}.
  \label{eq:rho}
\end{equation}
The queue is stable only when $\rho < 1$.
Throughout the thesis, I assume that arrivals follow a Poisson process and that service times are exponential, as in the standard treatment \cite{example2021}.
Chapter~\ref{ch:desks} applies this model to the return desk of a tool library.
